% !TEX program = latexmk
% !TEX options = -xelatex -synctex=1 -interaction=nonstopmode -file-line-error "%DOC%"
\documentclass[12pt, a4paper]{article}

\usepackage{fontspec}
\usepackage{xeCJK}
\xeCJKsetup{CJKmath=true, AutoFakeBold=2.5, AutoFakeSlant=0.2}
\punctstyle{quanjiao}

% ── 中文字体（使用系统可用的文泉驿正黑） ──
\IfFontExistsTF{WenQuanYi Zen Hei}{
  \setCJKmainfont{WenQuanYi Zen Hei}
  \setCJKsansfont{WenQuanYi Zen Hei}
  \IfFontExistsTF{WenQuanYi Zen Hei Mono}{
    \setCJKmonofont{WenQuanYi Zen Hei Mono}
  }{
    \setCJKmonofont{WenQuanYi Zen Hei}
  }
}{
  \IfFontExistsTF{Microsoft YaHei}{
    \setCJKmainfont{Microsoft YaHei}
    \setCJKsansfont{Microsoft YaHei}
    \setCJKmonofont{Microsoft YaHei}
  }{
    \setCJKmainfont{SimSun}
    \setCJKsansfont{SimSun}
    \setCJKmonofont{SimSun}
  }
}

% ── 西文字体 ──
\IfFontExistsTF{Liberation Serif}{
  \setmainfont{Liberation Serif}
}{
  \setmainfont{Times New Roman}
}
\IfFontExistsTF{Liberation Sans}{
  \setsansfont{Liberation Sans}
}{
  \setsansfont{Arial}
}
\IfFontExistsTF{DejaVu Sans Mono}{
  \setmonofont{DejaVu Sans Mono}
}{
  \setmonofont{Consolas}
}

\usepackage{geometry}
\geometry{
  a4paper,
  top=3cm, bottom=3cm,
  left=3.5cm, right=3.5cm,
  headheight=15pt,
  headsep=0.8cm,
  footskip=1.2cm
}

\usepackage{setspace}
\setstretch{1.5}

\usepackage{indentfirst}
\setlength{\parindent}{2em}
\setlength{\parskip}{0.4ex plus 0.2ex minus 0.1ex}

\usepackage[dvipsnames,svgnames,x11names]{xcolor}
\definecolor{primary}{RGB}{27, 79, 138}
\definecolor{secondary}{RGB}{46, 109, 173}
\definecolor{lightblue}{RGB}{235, 244, 255}
\definecolor{tabhead}{RGB}{27, 79, 138}
\definecolor{rowalt}{RGB}{240, 245, 251}

\usepackage{graphicx}
\usepackage{tikz}
\usetikzlibrary{shapes.geometric, arrows.meta, positioning, calc}

\usepackage{titlesec}
\titleformat{\section}{\large\bfseries\color{primary}}{\thesection}{1em}{}[{\color{secondary}\titlerule[0.8pt]}]
\titleformat{\subsection}{\normalsize\bfseries\color{secondary}}{\thesubsection}{1em}{}
\titleformat{\subsubsection}{\normalsize\color{secondary}}{\thesubsubsection}{1em}{}
\titlespacing{\section}{0pt}{20pt}{8pt}
\titlespacing{\subsection}{0pt}{14pt}{5pt}
\titlespacing{\subsubsection}{0pt}{10pt}{4pt}

\usepackage{fancyhdr}
\pagestyle{fancy}
\fancyhf{}
\fancyhead[L]{\small\color{primary}\bfseries 基于5G基站的无人机ISAR成像感知系统}
\fancyhead[R]{\small\color{gray} 技术研究报告}
\fancyfoot[C]{\small\color{gray}\thepage}
\renewcommand{\headrulewidth}{0.6pt}
\renewcommand{\headrule}{\color{primary}\hrule width\headwidth height\headrulewidth}

\usepackage{booktabs}
\usepackage{tabularx}
\usepackage{colortbl}
\usepackage{array}
\newcolumntype{C}{>{\centering\arraybackslash}X}
\newcolumntype{L}{>{\raggedright\arraybackslash}X}

\usepackage{enumitem}
\setlist[itemize]{leftmargin=2em, itemsep=3pt, topsep=5pt, parsep=0pt}
\setlist[enumerate]{leftmargin=2em, itemsep=3pt, topsep=5pt, parsep=0pt}

\usepackage{amsmath, amssymb, bm}
\usepackage{float}
% 公式后自动恢复首行缩进
\usepackage{etoolbox}
\AtBeginEnvironment{equation}{\setlength{\belowdisplayskip}{8pt plus 2pt minus 2pt}}
\usepackage[labelfont=bf, font=small, skip=6pt]{caption}
\usepackage[colorlinks=true, linkcolor=primary, urlcolor=secondary, citecolor=primary]{hyperref}

\usepackage{tcolorbox}
\tcbuselibrary{skins, breakable}

\newtcolorbox{infobox}{
  enhanced, breakable,
  colback=lightblue, colframe=primary,
  leftrule=4pt, rightrule=0pt, toprule=0.4pt, bottomrule=0.4pt,
  arc=2pt, left=10pt, right=10pt, top=8pt, bottom=8pt,
}

\newtcolorbox{metricbox}{
  enhanced,
  colback=primary!8, colframe=primary!50,
  arc=4pt, boxrule=0.6pt,
  left=12pt, right=12pt, top=10pt, bottom=10pt,
}

% 图片占位符命令：仿真结果图跑出来后替换即可
\newcommand{\figplaceholder}[2]{%
  \begin{tcolorbox}[enhanced, colback=gray!8, colframe=gray!40,
    arc=4pt, boxrule=0.5pt, left=10pt, right=10pt, top=15pt, bottom=15pt]
    \centering\color{gray}\small
    【待插入仿真结果图】\\[4pt] #1 \\[2pt]
    \texttt{#2}
  \end{tcolorbox}%
}

% 表格行距
\renewcommand{\arraystretch}{1.35}

\begin{document}
\thispagestyle{empty}

% ─── 封面 ───
\begin{tikzpicture}[remember picture, overlay]
  \fill[primary] (current page.north west)
    rectangle ([yshift=-5.8cm]current page.north east);
  \fill[secondary!70]
    ([yshift=-5.8cm]current page.north west)
    rectangle ([yshift=-6.5cm]current page.north east);
\end{tikzpicture}

\vspace*{1.2cm}
{\fontsize{30pt}{36pt}\selectfont\bfseries\color{white} 技术研究报告}\\[0.5cm]
{\large\color{white!80} Technical Research Report}

\vspace{4.6cm}

\begin{center}
{\fontsize{17pt}{23pt}\selectfont\bfseries\color{primary}
基于5G基站毫米波信号的无人机\\[0.25cm]
ISAR成像与智能感知系统}\\[0.6cm]
{\normalsize\color{gray} UAV Detection, ISAR Imaging and Tracking via 5G mmWave}\\[0.2cm]

\end{center}

\vspace{1.8cm}
\begin{center}
\renewcommand{\arraystretch}{1.6}
\begin{tabular}{rl}
\color{gray}所在单位   & 重庆邮电大学\quad 通信与信息工程学院 \\
\color{gray}研究方向   & 智能感知与雷达成像 \\
\color{gray}代码仓库   &
  \href{https://github.com/A1RER/Radar_Proj}{A1RER/Radar\_Proj} \\
\color{gray}报告日期   & 2026年3月 \\
\color{gray}版本       & v2.0 \\
\end{tabular}
\end{center}

\newpage
\setcounter{page}{1}

% ─── 摘要 ───
\begin{infobox}
{\bfseries\large\color{primary} 摘\quad 要}

\medskip
本项目以5G NR FR2毫米波基站信号（载频$f_c=28\,$GHz、带宽$B=400\,$MHz）为雷达照射源，构建了面向低空无人机（UAV）目标的检测、成像与跟踪一体化感知系统。系统核心由以下四大算法模块组成：

\medskip
\textbf{（1）ISAR成像}：利用线性调频（LFM）脉冲信号与目标旋转运动产生的多普勒频移，经距离压缩、运动补偿与自聚焦处理，生成二维高分辨率雷达图像。在400\,MHz带宽条件下，距离分辨率达\textbf{37.5\,cm}；在观测时间$T_{\mathrm{obs}}=2\,$s、旋转角速度$\omega=0.5\,$rad/s条件下，方位分辨率达\textbf{5.36\,mm}。

\textbf{（2）PSO自聚焦优化}：以最小化图像熵为目标函数，采用粒子群优化（PSO）算法自动搜索最优相位误差补偿系数，图像熵降低约\textbf{35\%}，20次迭代内稳定收敛。

\textbf{（3）卡尔曼滤波跟踪}：基于匀速（CV）运动模型建立四维状态方程，目标轨迹RMSE由$5.2\,$m降至\textbf{1.8\,m}，定位精度提升约\textbf{65\%}。

\textbf{（4）多基站协同定位}：实现TDOA/AOA/RSS三种定位体制及其加权融合算法，结合GDOP几何精度分析，为多5G基站协同探测提供理论评估框架。

\medskip
系统提供MATLAB与Python完整双语实现，并集成SVM、随机森林、CNN等机器学习分类子系统，支持5类无人机目标识别。全部实验基于仿真数据验证，后续将向实测数据与系统联调方向推进。

\medskip
\textbf{关键词：}5G通感一体化\quad ISAR成像\quad 粒子群优化\quad 卡尔曼滤波\quad 无人机检测\quad 毫米波雷达\quad 多基站定位
\end{infobox}

\vspace{0.6cm}
\tableofcontents
\newpage

% ════════════════════════════════════════════════════════
\section{研究背景与意义}

\subsection{低空无人机安全管控的现实需求}

随着消费级与工业级无人机技术的快速普及，低空空域安全管控问题日益突出。据统计，近年来"黑飞"事件频发，对机场净空区、军事敏感区域及城市公共安全构成严峻威胁\cite{ref_uav_threat}。现有低空目标探测手段主要包括光学摄像、声学探测和专用雷达系统，但分别存在受天气影响大、探测距离短、部署成本高等不足，迫切需要一种低成本、广覆盖、全天候的探测方案。

\subsection{5G通感一体化的技术机遇}

第五代移动通信（5G）NR标准定义了FR2频段（毫米波，24.25--52.6\,GHz），其中n257频段中心频率为28\,GHz，支持最大400\,MHz的信号带宽\cite{ref_3gpp}。毫米波信号具有以下两个对雷达感知至关重要的物理特性：

\begin{enumerate}
  \item \textbf{高距离分辨率}：雷达距离分辨率由信号带宽决定，$\Delta R = c/(2B)$。400\,MHz带宽对应$\Delta R = 37.5\,$cm，足以分辨无人机各主要散射部件；
  \item \textbf{短波长与强散射}：28\,GHz对应波长$\lambda = c/f_c \approx 1.07\,$cm，与无人机典型尺寸（0.2--2\,m）相比远小于目标尺度，处于光学区散射范围，目标散射截面积（RCS）较为稳定。
\end{enumerate}

利用现有5G基站基础设施兼做雷达感知节点，实现\textbf{通信-感知一体化}（Integrated Sensing and Communication, ISAC），可在不增加硬件部署的前提下大幅降低低空监控成本\cite{ref_isac_survey}。这一思路已成为6G关键技术方向之一。

\subsection{ISAR成像的适用性分析}

合成孔径雷达（SAR）通过雷达平台移动合成大孔径来获得方位向高分辨率，要求雷达载体精确运动。而\textbf{逆合成孔径雷达}（Inverse SAR, ISAR）则利用\textbf{目标自身}相对于雷达的旋转运动来合成等效大孔径\cite{ref_isar_chen}，其工作原理如图~\ref{fig:sar_vs_isar}所示。

\begin{figure}[H]
\centering
\begin{tikzpicture}[
  node distance=1.2cm,
  block/.style={draw=primary, fill=lightblue, rounded corners=3pt,
    minimum width=3.8cm, minimum height=0.7cm,
    font=\small, text=primary, align=center},
  label/.style={font=\small\color{gray}, align=center},
]
  % SAR部分
  \node[label] (sar_title) {\textbf{SAR（合成孔径雷达）}};
  \node[block, below=0.3cm of sar_title] (sar1) {雷达平台移动};
  \node[block, below=0.4cm of sar1] (sar2) {目标静止};
  \node[block, below=0.4cm of sar2] (sar3) {平台轨迹合成大孔径};

  % ISAR部分
  \node[label, right=3cm of sar_title] (isar_title) {\textbf{ISAR（逆合成孔径雷达）}};
  \node[block, below=0.3cm of isar_title, fill=primary!15] (isar1) {雷达固定（5G基站）};
  \node[block, below=0.4cm of isar1, fill=primary!15] (isar2) {目标旋转运动};
  \node[block, below=0.4cm of isar2, fill=primary!15] (isar3) {目标旋转合成大孔径};
\end{tikzpicture}
\caption{SAR与ISAR工作原理对比}
\label{fig:sar_vs_isar}
\end{figure}

无人机飞行过程中因姿态调整、风扰和机动操作而产生持续的旋转角度变化，这为ISAR成像提供了天然的多普勒历程。相比SAR，ISAR对雷达平台无运动要求，非常适合部署于固定的5G基站节点。

% ════════════════════════════════════════════════════════
\section{系统总体设计}

\subsection{系统架构}

系统按信号处理流程划分为五个层次，各层之间以矩阵形式传递数据，模块间完全解耦，支持独立调试与功能替换。整体架构如图~\ref{fig:system_arch}所示。

\begin{figure}[H]
\centering
\begin{tikzpicture}[
  node distance=0.65cm,
  box/.style={
    draw=primary, fill=lightblue, rounded corners=4pt,
    minimum width=6.2cm, minimum height=0.76cm,
    font=\small\bfseries, text=primary, align=center},
  arrow/.style={-Stealth, thick, color=secondary},
  side/.style={
    draw=secondary!60, fill=secondary!10, rounded corners=3pt,
    minimum width=3.2cm, minimum height=0.65cm,
    font=\small, align=center},
]
  \node[box] (sig)  {5G基站信号 {\small\mdseries $f_c$=28\,GHz，$B$=400\,MHz}};
  \node[box, below=of sig]  (echo) {回波数据生成 {\small\mdseries 4散射点，SNR=10\,dB}};
  \node[box, below=of echo] (rc)   {距离压缩（匹配滤波）};
  \node[box, below=of rc]   (mc)   {运动补偿（包络对齐 + 相位校正）};
  \node[box, below=of mc]   (af)   {自聚焦（PSO熵最小化）};
  \node[box, below=of af, fill=primary!18, draw=primary]
    (isar) {\bfseries ISAR 二维高分辨率图像};
  \node[side, right=1.5cm of mc] (ml) {ML分类\\目标识别};
  \node[side, right=1.5cm of isar] (pso) {PSO优化\\图像质量评估};
  \node[side, below=0.6cm of pso]  (kal) {卡尔曼跟踪\\轨迹滤波};
  \node[side, below=0.6cm of kal]  (loc) {多基站定位\\TDOA/AOA/RSS};
  \draw[arrow](sig)--(echo); \draw[arrow](echo)--(rc);
  \draw[arrow](rc)--(mc);    \draw[arrow](mc)--(af);
  \draw[arrow](af)--(isar);
  \draw[arrow](isar.east)--(pso.west);
  \draw[arrow](isar.east)--++(0.6,0)|-(kal.west);
  \draw[arrow](kal.south)|-(loc.west);
  \draw[arrow](isar.east)--++(0.6,0)|-(ml.west);
\end{tikzpicture}
\caption{系统信号处理流程图}
\label{fig:system_arch}
\end{figure}

\subsection{关键参数设置}

表~\ref{tab:params}列出了系统仿真所用的关键参数及其物理含义。参数选取依据5G NR FR2标准规范与典型低空无人机探测场景。

\begin{table}[H]
\centering
\begin{tabularx}{\textwidth}{lCCL}
\toprule
\rowcolor{tabhead}\color{white}\bfseries 参数 &
  \color{white}\bfseries 符号 &
  \color{white}\bfseries 取值 &
  \color{white}\bfseries 物理含义 \\
\midrule
载波频率 & $f_c$ & 28\,GHz & 5G NR n257频段中心频率 \\
\rowcolor{rowalt}
波长 & $\lambda$ & 1.07\,cm & $\lambda = c/f_c$ \\
信号带宽 & $B$ & 400\,MHz & 距离分辨率 $\Delta R = c/(2B) = 37.5\,$cm \\
\rowcolor{rowalt}
脉冲宽度 & $T_p$ & $1\,\mu$s & LFM脉冲持续时间 \\
调频率 & $K_r$ & $4\times10^{14}\,$Hz/s & $K_r = B/T_p$ \\
\rowcolor{rowalt}
脉冲重复频率 & PRF & 1000\,Hz & 方位向采样率 \\
观测时间 & $T_{\mathrm{obs}}$ & 2\,s & 相干积累时间 \\
\rowcolor{rowalt}
目标旋转角速度 & $\omega$ & 0.5\,rad/s & 等效转台角速度 \\
初始目标距离 & $R_0$ & 1000\,m & 雷达至目标中心距离 \\
\rowcolor{rowalt}
散射点模型 & --- & 4点十字形 & 模拟典型四旋翼结构 \\
信噪比 & SNR & 10\,dB & 模拟实际信道环境 \\
\rowcolor{rowalt}
卡尔曼采样间隔 & $\Delta t$ & 0.1\,s & 匀速运动模型步长 \\
\bottomrule
\end{tabularx}
\caption{系统关键参数}
\label{tab:params}
\end{table}

% ════════════════════════════════════════════════════════
\section{信号模型与理论基础}
\label{sec:signal_model}

本节从发射信号模型出发，逐步推导回波信号表达式、距离压缩原理与分辨率公式，为后续成像处理提供完整的数学基础。

\subsection{线性调频（LFM）发射信号模型}

系统采用线性调频（Linear Frequency Modulation, LFM）脉冲信号作为探测波形。LFM信号的瞬时频率随时间线性增长，能够在有限的发射功率下获得较大的时间-带宽积，从而通过脉冲压缩技术实现高距离分辨率。

单个LFM脉冲的基带表达式为：
\begin{equation}
  s_t(\hat{t}) = \mathrm{rect}\!\left(\frac{\hat{t}}{T_p}\right)
  \exp\!\left(j\pi K_r \hat{t}^{\,2}\right)
  \label{eq:lfm_baseband}
\end{equation}
其中，$\hat{t}$为快时间（脉冲内时间），$T_p$为脉冲宽度，$K_r = B/T_p$为调频率（单位Hz/s），$\mathrm{rect}(\cdot)$为矩形窗函数。

经射频上变频后的发射信号为：
\begin{equation}
  s_{\mathrm{RF}}(\hat{t}) = \mathrm{rect}\!\left(\frac{\hat{t}}{T_p}\right)
  \exp\!\left(j2\pi f_c \hat{t} + j\pi K_r \hat{t}^{\,2}\right)
  \label{eq:lfm_rf}
\end{equation}

LFM信号的瞬时频率为：
\begin{equation}
  f_i(\hat{t}) = f_c + K_r \hat{t}
  \label{eq:inst_freq}
\end{equation}
在脉冲持续时间$T_p$内，频率从$f_c - B/2$线性扫过至$f_c + B/2$，总带宽为$B$。

\subsection{无人机散射点目标模型}

将无人机建模为$K$个等效散射点的集合体。在雷达坐标系下，第$k$个散射点的初始坐标为$(x_k, y_k)$，散射系数为$\sigma_k$。目标以角速度$\omega$绕几何中心旋转（模拟飞行姿态变化），则在慢时间$t_m$（第$m$个脉冲对应的时刻）时刻，第$k$个散射点的旋转后坐标为：
\begin{equation}
  \begin{cases}
    x_k'(t_m) = x_k \cos(\omega t_m) - y_k \sin(\omega t_m) \\[4pt]
    y_k'(t_m) = x_k \sin(\omega t_m) + y_k \cos(\omega t_m)
  \end{cases}
  \label{eq:rotation}
\end{equation}

本项目采用4散射点十字形模型来模拟典型的四旋翼无人机结构，散射点分布如表~\ref{tab:scatter}所示。

\begin{table}[H]
\centering
\begin{tabularx}{0.7\textwidth}{cCCC}
\toprule
\rowcolor{tabhead}\color{white}\bfseries 散射点 &
  \color{white}\bfseries $x_k$ (m) &
  \color{white}\bfseries $y_k$ (m) &
  \color{white}\bfseries 物理含义 \\
\midrule
1 & 0 & 0 & 机体中心 \\
\rowcolor{rowalt}
2 & 0.3 & 0 & 右侧机臂 \\
3 & $-$0.3 & 0 & 左侧机臂 \\
\rowcolor{rowalt}
4 & 0 & 0.25 & 前方机臂 \\
\bottomrule
\end{tabularx}
\caption{无人机散射点模型参数}
\label{tab:scatter}
\end{table}

\subsection{回波信号模型}

第$k$个散射点在第$m$个脉冲时刻的瞬时距离为：
\begin{equation}
  R_k(t_m) = R_0 + y_k'(t_m)
  \label{eq:range}
\end{equation}
其中$R_0$为目标中心到雷达的基准距离。信号的双程传播时延为$\tau_k = 2R_k/c$。

总回波信号为所有散射点回波的叠加，第$m$个脉冲的回波为：
\begin{equation}
  s_r(\hat{t},\, t_m) = \sum_{k=1}^{K} \sigma_k \cdot
  \mathrm{rect}\!\left(\frac{\hat{t}-\tau_k}{T_p}\right)
  \exp\!\left[j\pi K_r\!\left(\hat{t}-\tau_k\right)^{2}\right]
  \exp\!\left(-j\frac{4\pi f_c R_k(t_m)}{c}\right)
  + n(\hat{t},\, t_m)
  \label{eq:echo}
\end{equation}
其中，第一个指数项为LFM脉冲压缩相位，第二个指数项包含了距离与多普勒信息，$n(\hat{t},\allowbreak t_m)$为加性高斯白噪声。

\subsection{匹配滤波与距离压缩}
\label{subsec:matched_filter}

\textbf{匹配滤波}是最优线性检测方法。其原理是：用发射信号的共轭时间反转作为滤波器，使输出信噪比最大化\cite{ref_radar_signal}。

匹配滤波器的冲激响应为：
\begin{equation}
  h(\hat{t}) = s_t^{*}(-\hat{t}) = \exp\!\left(-j\pi K_r \hat{t}^{\,2}\right)
  \label{eq:matched_filter}
\end{equation}

回波信号与匹配滤波器做卷积后，输出为压缩脉冲：
\begin{equation}
  s_c(\hat{t}) = s_r(\hat{t}) * h(\hat{t})
  \approx B T_p \cdot \mathrm{sinc}\!\left[B\!\left(\hat{t} - \tau_k\right)\right]
  \exp\!\left(-j\frac{4\pi f_c R_k}{c}\right)
  \label{eq:pulse_compression}
\end{equation}

sinc函数的主瓣宽度决定了距离分辨率。$\mathrm{sinc}$函数的3\,dB主瓣宽度约为$1/B$（时间），对应的距离分辨率为：
\begin{equation}
  \boxed{\Delta R = \frac{c}{2B} = \frac{3\times 10^8}{2\times 400\times 10^6} = 0.375\;\text{m} = 37.5\;\text{cm}}
  \label{eq:range_res}
\end{equation}

脉冲压缩的物理意义在于：通过匹配滤波将时间宽度为$T_p$的长脉冲压缩为宽度约$1/B$的窄脉冲，在不增加峰值发射功率的条件下实现了高距离分辨率。\textbf{时间-带宽积}$BT_p$衡量了压缩增益：
\begin{equation}
  BT_p = 400\times 10^6 \times 1\times 10^{-6} = 400
  \label{eq:tbp}
\end{equation}

\subsection{方位向多普勒分辨与ISAR成像原理}
\label{subsec:azimuth}

距离压缩完成了快时间维度（距离向）的处理。为获得方位向（慢时间维度）的分辨能力，需要利用目标旋转产生的多普勒频移。

第$k$个散射点的瞬时径向速度为：
\begin{equation}
  v_{r,k}(t_m) = \frac{dR_k}{dt_m} = \omega\!\left[x_k\cos(\omega t_m) + y_k\sin(\omega t_m)\right] \approx \omega\, x_k
  \label{eq:radial_velocity}
\end{equation}
在$\omega T_{\mathrm{obs}} \ll 1$的小角度近似下，$\cos(\omega t_m) \approx 1$，$\sin(\omega t_m) \approx \omega t_m$。对应的多普勒频率为：
\begin{equation}
  f_{d,k} = \frac{2 v_{r,k}}{\lambda} = \frac{2\omega\, x_k}{\lambda}
  = \frac{2\omega\, x_k\, f_c}{c}
  \label{eq:doppler}
\end{equation}

由式~(\ref{eq:doppler})可知，不同方位位置$x_k$的散射点对应不同的多普勒频率，这就是ISAR方位向分辨的物理基础。对距离压缩后的慢时间数据做傅里叶变换（方位向FFT），即可在多普勒频率轴上分辨不同方位位置的散射点，从而形成二维ISAR图像。

\textbf{方位分辨率}由相干积累角度$\Delta\theta$决定：
\begin{equation}
  \Delta\theta = \omega \cdot T_{\mathrm{obs}}
  \label{eq:coherent_angle}
\end{equation}
\begin{equation}
  \boxed{\rho_a = \frac{\lambda}{2\Delta\theta} = \frac{\lambda}{2\omega T_{\mathrm{obs}}}
  = \frac{1.07\times 10^{-2}}{2\times 0.5\times 2} = 5.36\times 10^{-3}\;\text{m} \approx 5.36\;\text{mm}}
  \label{eq:azimuth_res}
\end{equation}

需要指出，上述方位分辨率为理论极限值，实际成像中受运动补偿精度、相位误差等因素影响，真实分辨率通常低于该值。

\subsection{ISAR成像完整处理流程}

综合以上理论推导，ISAR成像的完整处理步骤如下：

\begin{enumerate}
  \item \textbf{距离压缩}：对每个慢时间采样（脉冲）的回波信号与匹配滤波器卷积，将展宽的回波压缩为窄脉冲，得到距离像。这一步利用式~(\ref{eq:matched_filter})--(\ref{eq:pulse_compression})实现；
  \item \textbf{包络对齐（平动补偿）}：目标的平动（整体径向运动）导致各脉冲的距离像发生平移。采用\textbf{相关法}——以某一参考脉冲的距离像为模板，通过互相关找出其余脉冲的平移量并予以校正；
  \item \textbf{相位校正（自聚焦）}：包络对齐后仍存在残余相位误差，需通过自聚焦算法（如最小熵法或PGA）进一步补偿。本项目采用PSO全局优化进行自聚焦（详见第~\ref{sec:pso}节）；
  \item \textbf{方位向FFT}：对补偿后的二维数据沿慢时间轴做傅里叶变换，利用式~(\ref{eq:doppler})的多普勒-方位映射关系，输出最终的二维ISAR图像。
\end{enumerate}

% ════════════════════════════════════════════════════════
\section{核心算法设计}

\subsection{PSO图像质量优化（自聚焦）}
\label{sec:pso}

\subsubsection{问题建模}

运动补偿后的残余相位误差可建模为慢时间的多项式函数：
\begin{equation}
  \phi_{\mathrm{err}}(m) = \alpha\, m + \beta\, m^2
  \label{eq:phase_error}
\end{equation}
其中$m$为脉冲索引，$\alpha$和$\beta$分别为一次和二次相位误差系数。校正后的信号为：
\begin{equation}
  s_{\mathrm{corr}}(\hat{t},\, m) = s_c(\hat{t},\, m) \cdot \exp\!\left[-j\phi_{\mathrm{err}}(m)\right]
  \label{eq:phase_correction}
\end{equation}

\subsubsection{图像熵目标函数}

\textbf{信息熵}是衡量ISAR图像聚焦质量的经典指标\cite{ref_isar_autofocus}。对于一幅ISAR图像$I(i,j)$，定义归一化强度分布：
\begin{equation}
  p(i,j) = \frac{|I(i,j)|}{\displaystyle\sum_{i,j}|I(i,j)|}
  \label{eq:entropy_dist}
\end{equation}

图像熵定义为：
\begin{equation}
  H(I) = -\sum_{i,j} p(i,j)\,\ln\, p(i,j)
  \label{eq:entropy}
\end{equation}

聚焦良好的图像能量集中于少数像素，熵值较低；散焦图像能量分散，熵值较高。因此，自聚焦可转化为如下优化问题：
\begin{equation}
  \min_{\alpha,\,\beta}\; H\!\left(I(\alpha,\beta)\right)
  \quad\text{s.t.}\quad
  \alpha\in[-0.1,\;0.1],\quad
  \beta\in[-0.01,\;0.01]
  \label{eq:pso_objective}
\end{equation}

\subsubsection{粒子群优化（PSO）算法}

PSO是由Kennedy和Eberhart于1995年提出的群体智能优化算法\cite{ref_pso}，灵感来源于鸟群觅食行为。算法维护一组"粒子"在搜索空间中运动，每个粒子代表一个候选解。

设第$i$个粒子在第$t$次迭代时的位置为$\bm{x}_i^t$，速度为$\bm{v}_i^t$，则更新公式为：
\begin{equation}
  \bm{v}_i^{t+1} = w\,\bm{v}_i^t
  + c_1 r_1 \left(\bm{p}_i - \bm{x}_i^t\right)
  + c_2 r_2 \left(\bm{g} - \bm{x}_i^t\right)
  \label{eq:pso_velocity}
\end{equation}
\begin{equation}
  \bm{x}_i^{t+1} = \bm{x}_i^t + \bm{v}_i^{t+1}
  \label{eq:pso_position}
\end{equation}
其中：
\begin{itemize}
  \item $w = 0.7$ 为\textbf{惯性权重}，控制粒子对历史速度的继承程度；
  \item $c_1 = 1.5$ 为\textbf{个体学习因子}，引导粒子向自身历史最优位置$\bm{p}_i$靠近；
  \item $c_2 = 1.5$ 为\textbf{社会学习因子}，引导粒子向全局最优位置$\bm{g}$靠近；
  \item $r_1, r_2 \in [0, 1]$ 为均匀随机数，增加搜索的随机性。
\end{itemize}

相比梯度下降法，PSO具有以下优势：（1）不需要计算目标函数的梯度；（2）全局搜索能力强，能有效规避局部极值；（3）算法简单，易于实现。

\begin{table}[H]
\centering
\begin{tabularx}{0.72\textwidth}{LCC}
\toprule
\rowcolor{tabhead}\color{white}\bfseries PSO参数 &
  \color{white}\bfseries 符号 &
  \color{white}\bfseries 取值 \\
\midrule
粒子数 & $N_p$ & 20 \\
\rowcolor{rowalt}
搜索维度 & $D$ & 2（$\alpha, \beta$） \\
惯性权重 & $w$ & 0.7 \\
\rowcolor{rowalt}
个体学习因子 & $c_1$ & 1.5 \\
社会学习因子 & $c_2$ & 1.5 \\
\rowcolor{rowalt}
最大迭代次数 & $T_{\max}$ & 30 \\
\bottomrule
\end{tabularx}
\caption{PSO算法参数设置}
\label{tab:pso}
\end{table}

\subsection{卡尔曼滤波目标跟踪}
\label{sec:kalman}

\subsubsection{匀速运动模型（CV模型）}

假设无人机在短时间内做近似匀速直线运动，定义四维状态向量：
\begin{equation}
  \bm{x}_k = \begin{bmatrix} x_k \\ y_k \\ v_{x,k} \\ v_{y,k} \end{bmatrix}
  \label{eq:state_vector}
\end{equation}
其中$(x_k, y_k)$为平面位置，$(v_{x,k}, v_{y,k})$为速度分量。

\textbf{状态转移方程}（预测模型）为：
\begin{equation}
  \bm{x}_k = \bm{F}\,\bm{x}_{k-1} + \bm{w}_{k-1}
  \label{eq:state_transition}
\end{equation}
其中状态转移矩阵为：
\begin{equation}
  \bm{F} = \begin{bmatrix}
    1 & 0 & \Delta t & 0 \\
    0 & 1 & 0 & \Delta t \\
    0 & 0 & 1 & 0 \\
    0 & 0 & 0 & 1
  \end{bmatrix}
  \label{eq:F_matrix}
\end{equation}
$\bm{w}_{k-1} \sim \mathcal{N}(\bm{0}, \bm{Q})$为过程噪声，$\bm{Q}$为过程噪声协方差矩阵，反映模型不确定性。

\textbf{观测方程}为：
\begin{equation}
  \bm{z}_k = \bm{H}\,\bm{x}_k + \bm{v}_k,\qquad
  \bm{H} = \begin{bmatrix} 1 & 0 & 0 & 0 \\ 0 & 1 & 0 & 0 \end{bmatrix}
  \label{eq:observation}
\end{equation}
$\bm{v}_k \sim \mathcal{N}(\bm{0}, \bm{R})$为测量噪声，$\bm{R}$为测量噪声协方差矩阵。

\subsubsection{卡尔曼滤波递推公式}

卡尔曼滤波由\textbf{预测}和\textbf{更新}两步交替进行：

\textbf{预测步骤：}
\begin{equation}
  \hat{\bm{x}}_{k|k-1} = \bm{F}\,\hat{\bm{x}}_{k-1|k-1}
  \label{eq:kf_predict_state}
\end{equation}
\begin{equation}
  \bm{P}_{k|k-1} = \bm{F}\,\bm{P}_{k-1|k-1}\,\bm{F}^\top + \bm{Q}
  \label{eq:kf_predict_cov}
\end{equation}

\textbf{更新步骤：}
\begin{equation}
  \bm{K}_k = \bm{P}_{k|k-1}\,\bm{H}^\top
  \left(\bm{H}\,\bm{P}_{k|k-1}\,\bm{H}^\top + \bm{R}\right)^{-1}
  \label{eq:kf_gain}
\end{equation}
\begin{equation}
  \hat{\bm{x}}_{k|k} = \hat{\bm{x}}_{k|k-1}
  + \bm{K}_k\!\left(\bm{z}_k - \bm{H}\,\hat{\bm{x}}_{k|k-1}\right)
  \label{eq:kf_update_state}
\end{equation}
\begin{equation}
  \bm{P}_{k|k} = \left(\bm{I} - \bm{K}_k\,\bm{H}\right)\bm{P}_{k|k-1}
  \label{eq:kf_update_cov}
\end{equation}

其中，$\bm{K}_k$为\textbf{卡尔曼增益}矩阵，自动权衡预测值与观测值的可信度：当测量噪声$\bm{R}$较大时，$\bm{K}_k$较小，滤波器更信任预测；反之更信任观测。$\bm{P}_{k|k}$为估计误差协方差矩阵，反映当前状态估计的不确定性。

\textbf{跟踪性能评估}采用均方根误差（RMSE）：
\begin{equation}
  \mathrm{RMSE} = \sqrt{\frac{1}{N}\sum_{k=1}^{N}
  \left[(x_k - \hat{x}_k)^2 + (y_k - \hat{y}_k)^2\right]}
  \label{eq:rmse}
\end{equation}

\begin{metricbox}
\centering
{\bfseries 跟踪性能}\\[6pt]
原始测量 RMSE：5.2\,m
\quad$\xrightarrow{\text{卡尔曼滤波}}$\quad
\textbf{RMSE：1.8\,m}（精度提升约 \textbf{65\%}）
\end{metricbox}

\subsection{多基站协同定位}
\label{sec:multistation}

当多个5G基站同时观测同一目标时，可利用各基站的测量信息进行协同定位，显著提升空间定位精度。本项目实现了三种经典定位体制。

\subsubsection{TDOA定位（到达时间差）}

设$N$个基站的位置为$\bm{s}_i$（$i=1,\ldots,N$），目标位置为$\bm{p}$，信号从目标到第$i$个基站的传播距离为$d_i = \|\bm{p} - \bm{s}_i\|$。以第1个基站为参考，到达时间差测量为：
\begin{equation}
  \tau_{i1} = \frac{d_i - d_1}{c},\quad i = 2,\ldots,N
  \label{eq:tdoa}
\end{equation}

每个时间差测量定义一个以两个基站为焦点的双曲面，目标位置为所有双曲面的交点。将问题线性化，可得到最小二乘估计：
\begin{equation}
  \bm{A}\,\bm{p} = \bm{b}
  \label{eq:tdoa_linear}
\end{equation}
其中$\bm{A}$和$\bm{b}$由基站位置与时差测量构建。初始解经非线性优化精化后得到最终估计\cite{ref_tdoa}。

\subsubsection{AOA定位（到达角度）}

每个基站测量目标的方位角$\theta_i$和俯仰角$\phi_i$，形成一条从基站出发的方向射线：
\begin{equation}
  \bm{d}_i = \begin{bmatrix}
    \cos\phi_i\cos\theta_i \\
    \cos\phi_i\sin\theta_i \\
    \sin\phi_i
  \end{bmatrix}
  \label{eq:aoa_direction}
\end{equation}

目标位置为所有射线的最近交点，通过最小化点到各射线距离平方和求解：
\begin{equation}
  \hat{\bm{p}} = \arg\min_{\bm{p}} \sum_{i=1}^{N}
  \left\|\left(\bm{p}-\bm{s}_i\right) - \left[\left(\bm{p}-\bm{s}_i\right)^\top\bm{d}_i\right]\bm{d}_i\right\|^2
  \label{eq:aoa_objective}
\end{equation}

定义\textbf{垂直投影矩阵}$\bm{M}_i = \bm{I} - \bm{d}_i\bm{d}_i^\top$，它将任意向量投影到射线$\bm{d}_i$的正交补空间。由于$\bm{M}_i$对称且幂等（$\bm{M}_i^\top = \bm{M}_i$，$\bm{M}_i^2 = \bm{M}_i$），式~(\ref{eq:aoa_objective})可改写为：
\begin{equation}
  J(\bm{p}) = \sum_{i=1}^{N} (\bm{p} - \bm{s}_i)^\top \bm{M}_i (\bm{p} - \bm{s}_i)
  \label{eq:aoa_quadratic}
\end{equation}
这是关于$\bm{p}$的二次函数。对$\bm{p}$求导并令其为零：
\[
  \frac{\partial J}{\partial \bm{p}} = 2\sum_{i=1}^{N} \bm{M}_i(\bm{p} - \bm{s}_i) = \bm{0}
\]
整理得到线性方程$\bm{A}\bm{p} = \bm{b}$，其中：
\begin{equation}
  \bm{A} = \sum_{i=1}^{N}\bm{M}_i, \quad \bm{b} = \sum_{i=1}^{N}\bm{M}_i\bm{s}_i
  \label{eq:aoa_Ab}
\end{equation}
因此AOA定位存在\textbf{解析解}：
\begin{equation}
  \boxed{\hat{\bm{p}} = \bm{A}^{-1}\bm{b} = \left(\sum_{i=1}^{N}\bm{M}_i\right)^{-1}\left(\sum_{i=1}^{N}\bm{M}_i\bm{s}_i\right)}
  \label{eq:aoa_closed_form}
\end{equation}

投影矩阵$\bm{M}_i$仅由已知的方向向量$\bm{d}_i$决定，不依赖于待求量$\bm{p}$，因此目标函数是$\bm{p}$的二次型，令导数为零即可得到唯一的闭合解，无需数值迭代。这本质上是\textbf{最小二乘直线交汇}（Least-Squares Line Intersection）问题。相比TDOA定位中必须的非线性优化精化步骤，AOA解析解的计算开销极低，适合实时系统部署。

\subsubsection{RSS定位（接收信号强度）}

RSS定位基于无线信号的路径损耗模型。第$i$个基站接收到的信号强度为：
\begin{equation}
  P_{r,i} = P_0 - 10n\log_{10}\!\left(\frac{d_i}{d_0}\right) + \xi_i
  \label{eq:rss_model}
\end{equation}
其中$P_0$为参考距离$d_0$处的接收功率，$n$为路径损耗指数（自由空间$n=2$），$\xi_i$为阴影衰落（对数正态分布）。由RSS值反算距离后，利用三边定位法（trilateration）求解目标位置。

\subsubsection{GDOP几何精度分析}

\textbf{几何精度衰减因子}（Geometric Dilution of Precision, GDOP）衡量基站几何构型对定位精度的影响。定义几何矩阵：
\begin{equation}
  \bm{G} = \begin{bmatrix}
    \dfrac{\partial d_1}{\partial x} & \dfrac{\partial d_1}{\partial y} & \dfrac{\partial d_1}{\partial z} & 1 \\[6pt]
    \vdots & \vdots & \vdots & \vdots \\[2pt]
    \dfrac{\partial d_N}{\partial x} & \dfrac{\partial d_N}{\partial y} & \dfrac{\partial d_N}{\partial z} & 1
  \end{bmatrix}
  \label{eq:gdop_G}
\end{equation}

GDOP定义为：
\begin{equation}
  \mathrm{GDOP} = \sqrt{\mathrm{tr}\!\left[\left(\bm{G}^\top\bm{G}\right)^{-1}\right]}
  \label{eq:gdop}
\end{equation}

GDOP值越小，基站构型越有利于定位。一般而言，GDOP $<$ 3为良好构型，GDOP $>$ 6时定位精度显著恶化。多基站协同的目标之一是通过合理布站使GDOP最小化。

% ════════════════════════════════════════════════════════
\section{无人机目标分类子系统}

\subsection{特征提取方法}

从ISAR图像中提取三类互补特征并拼接形成高维特征向量，送入分类器进行目标识别。

\subsubsection{HOG特征（方向梯度直方图）}

HOG（Histogram of Oriented Gradients）通过计算图像局部区域的梯度方向分布来描述目标形状\cite{ref_hog}。首先将图像划分为$8\times8$像素的cell，统计每个cell内各方向的梯度幅值直方图，再将相邻cell组成block进行归一化，最终拼接为一维特征向量。HOG特征对目标的外轮廓信息捕捉能力较强。

\subsubsection{LBP特征（局部二值模式）}

LBP（Local Binary Pattern）描述图像的局部纹理结构。对每个像素，将其与周围邻域像素进行大小比较，生成一个二进制编码，统计编码的直方图作为特征。LBP能够有效描述ISAR图像中散射点分布的纹理规律。

\subsubsection{统计特征}

从ISAR图像中提取全局统计量，包括：均值、标准差、最大值、最小值、方差，以及灰度共生矩阵（GLCM）导出的对比度、相关性、能量和均匀性等纹理特征。统计特征计算速度快，适合作为其他特征的补充。

\subsection{分类器设计与对比}

\begin{table}[H]
\centering
\begin{tabularx}{\textwidth}{lLLL}
\toprule
\rowcolor{tabhead}\color{white}\bfseries 分类器 &
  \color{white}\bfseries 核心原理 &
  \color{white}\bfseries 优势 &
  \color{white}\bfseries 适用场景 \\
\midrule
SVM（RBF核） & 寻找最大间隔超平面，RBF核映射到高维空间 & 泛化性强，适合中小样本 & 快速原型验证 \\
\rowcolor{rowalt}
随机森林 & 多棵决策树投票集成 & 可解释性好，不易过拟合 & 特征重要性分析 \\
KNN & 基于距离的近邻投票 & 无需训练，实现简单 & 基线对比 \\
\rowcolor{rowalt}
PyTorch CNN & 卷积层自动提取空间特征 & 端到端学习，精度高 & 大样本场景 \\
ResNet迁移学习 & 利用ImageNet预训练权重 & 数据需求少，收敛快 & 数据稀缺场景 \\
\bottomrule
\end{tabularx}
\caption{分类器对比}
\label{tab:classifiers}
\end{table}

\subsection{实验设置与评估指标}

仿真数据集包含1000张$128\times128$像素的合成ISAR图像，涵盖5类无人机目标：小型多旋翼、中型多旋翼、大型多旋翼、固定翼和旋翼机。数据集按70\%/30\%划分为训练集与测试集，并采用5折交叉验证评估泛化能力。

评估指标包括：
\begin{itemize}
  \item \textbf{准确率}（Accuracy）：正确分类样本占总样本的比例；
  \item \textbf{精确率}（Precision）：对每个类别，预测为该类的样本中实际属于该类的比例；
  \item \textbf{召回率}（Recall）：对每个类别，实际属于该类的样本中被正确预测的比例；
  \item \textbf{F1分数}：精确率与召回率的调和平均，$F_1 = 2 \cdot \frac{P \cdot R}{P + R}$。
\end{itemize}

% ════════════════════════════════════════════════════════
\section{仿真实验与结果分析}
\label{sec:experiments}

\subsection{ISAR成像结果}

\begin{figure}[H]
\centering
\includegraphics[width=0.95\textwidth]{../matlab_code/figures/fig_isar_result.png}
\caption{ISAR成像结果：(a) 原始回波数据；(b) 距离压缩后（dB显示）；(c) 二维ISAR图像（dB显示）}
\label{fig:isar_result}
\end{figure}

\subsection{PSO优化收敛过程}

\begin{figure}[H]
\centering
\includegraphics[width=0.65\textwidth]{../matlab_code/figures/fig_pso_convergence.png}
\caption{PSO优化收敛曲线：图像熵随迭代次数单调下降，约20次迭代后趋于稳定}
\label{fig:pso_convergence}
\end{figure}

\subsection{卡尔曼滤波跟踪结果}

\begin{figure}[H]
\centering
\includegraphics[width=0.95\textwidth]{../matlab_code/figures/fig_kalman_tracking.png}
\caption{卡尔曼滤波跟踪结果：(a) 真实轨迹、含噪测量与滤波估计对比；(b) RMSE定位精度对比}
\label{fig:kalman_tracking}
\end{figure}

\subsection{分类器性能对比}

\begin{figure}[H]
\centering
\includegraphics[width=0.6\textwidth]{../matlab_code/figures/fig_confusion_matrix.png}
\caption{无人机目标分类混淆矩阵（5类，基于统计特征+SVM分类器）}
\label{fig:confusion_matrix}
\end{figure}

% ════════════════════════════════════════════════════════
\section{系统实现}

\subsection{代码架构}

项目提供完整的MATLAB与Python双语实现，两者算法逻辑完全对应，便于在不同环境下交叉验证。代码结构如表~\ref{tab:code}所示。

\begin{table}[H]
\centering
\begin{tabularx}{\textwidth}{lLL}
\toprule
\rowcolor{tabhead}\color{white}\bfseries 功能模块 &
  \color{white}\bfseries MATLAB实现 &
  \color{white}\bfseries Python实现 \\
\midrule
ISAR成像 & \texttt{ISAR\_Demo.m} & \texttt{11\_ISAR\_Python.py} \\
\rowcolor{rowalt}
ISAR系统类 & \texttt{ISARImagingSystem.m} & \texttt{15\_ISAR\_Advanced.py} \\
PSO优化 & \texttt{PSO\_Optimization.m} & \texttt{12\_PSO.py} \\
\rowcolor{rowalt}
卡尔曼跟踪 & \texttt{Kalman\_Filter.m} & \texttt{13\_Kalman.py} \\
ML分类 & \texttt{MLClassifier.m} & \texttt{18\_ML\_Classification.py} \\
\rowcolor{rowalt}
SVM分类 & \texttt{Demo\_ML\_MATLAB.m} & \texttt{19\_SVM\_Classification.py} \\
合成数据生成 & \texttt{generateSyntheticData.m} & --- \\
\rowcolor{rowalt}
系统集成 & \texttt{Main\_System.m} & \texttt{16\_Main\_System.py} \\
多基站定位 & \texttt{\small MultiStationLocalization.m} & \texttt{\small 14\_MultiStation\_Localization.py} \\
\bottomrule
\end{tabularx}
\caption{双语代码对应关系}
\label{tab:code}
\end{table}

\subsection{运行环境与依赖}

\begin{itemize}
  \item \textbf{MATLAB}：R2020b及以上，需Signal Processing Toolbox、Optimization Toolbox（无工具箱时自动回退至手动实现）；
  \item \textbf{Python}：3.10+，主要依赖：\texttt{numpy}、\texttt{scipy}、\texttt{matplotlib}、\texttt{pyswarm}、\texttt{filterpy}、\texttt{torch}、\texttt{scikit-learn}。
\end{itemize}

% ════════════════════════════════════════════════════════
\section{性能指标汇总}

\begin{table}[H]
\centering
\begin{tabularx}{\textwidth}{lCCL}
\toprule
\rowcolor{tabhead}\color{white}\bfseries 性能指标 &
  \color{white}\bfseries 数值 &
  \color{white}\bfseries 对比基线 &
  \color{white}\bfseries 决定因素 \\
\midrule
距离分辨率 $\Delta R$ & 37.5\,cm & --- & $c/(2B)$，由带宽$B$决定 \\
\rowcolor{rowalt}
方位分辨率 $\rho_a$ & 5.36\,mm & --- & $\lambda/(2\omega T_\mathrm{obs})$ \\
时间-带宽积 $BT_p$ & 400 & --- & 脉冲压缩增益 \\
\rowcolor{rowalt}
PSO图像熵改善 & $\approx 35\%$ & 未优化直接成像 & 自聚焦补偿精度 \\
PSO收敛速度 & $<20$ 次迭代 & 梯度法约需50+次 & 全局搜索能力 \\
\rowcolor{rowalt}
卡尔曼跟踪RMSE & 1.8\,m & 原始测量5.2\,m & 滤波增益 \\
定位精度提升 & $\approx 65\%$ & 无滤波直接测量 & CV模型+噪声统计 \\
\rowcolor{rowalt}
支持目标类别数 & 5类 & --- & HOG+LBP+统计特征 \\
\bottomrule
\end{tabularx}
\caption{核心性能指标汇总}
\label{tab:performance}
\end{table}

% ════════════════════════════════════════════════════════
\section{现有不足}

本项目在仿真层面完成了系统性验证，但仍存在以下主要局限：

\begin{enumerate}
  \item \textbf{仿真数据驱动}：所有实验基于合成回波数据，尚未接入真实5G基站射频信号，无法验证实际信道条件下的算法鲁棒性；
  \item \textbf{目标模型简化}：4散射点旋转体模型对复杂飞行姿态（俯仰、偏航耦合）的表征精度有限，真实无人机的电磁散射特性远更复杂；
  \item \textbf{多基站融合待验}：协同定位模块已完成算法实现，但系统级多基站联合调试尚未开展，GDOP分析也需结合实际基站部署场景进行评估；
  \item \textbf{分类数据集较小}：1000张合成图像训练的分类器泛化能力有限，需扩大数据规模并引入数据增强策略；
  \item \textbf{方位分辨率依赖旋转}：ISAR方位分辨率取决于目标旋转角速度，当目标运动平稳（$\omega$接近0）时成像质量显著退化。
\end{enumerate}

% ════════════════════════════════════════════════════════
\section{结语}

本项目以5G NR FR2毫米波信号为切入点，系统性地构建了无人机ISAR成像感知的完整算法框架。从LFM信号模型与匹配滤波的理论推导出发，经距离压缩、运动补偿、PSO自聚焦优化，实现了高分辨率ISAR二维成像；在此基础上，结合卡尔曼滤波跟踪、多基站协同定位与机器学习目标分类，构建了"检测--成像--跟踪--识别"四位一体的感知系统。系统提供MATLAB与Python完整双语实现，在仿真层面验证了各核心算法的有效性。

本项目为5G通感一体化低空安防应用提供了完整的算法框架与仿真验证基础。

% ════════════════════════════════════════════════════════
\newpage
\begin{thebibliography}{99}

\bibitem{ref_uav_threat}
Shi X, Yang C, Xie W, et al.
Anti-drone system with multiple surveillance technologies: Architecture, implementation, and challenges[J].
\textit{IEEE Communications Magazine}, 2018, 56(4): 68--74.

\bibitem{ref_3gpp}
3GPP.
NR; Physical channels and modulation (Release 16)[S].
TS 38.211, 2020.

\bibitem{ref_isac_survey}
Liu F, Cui Y, Masouros C, et al.
Integrated sensing and communications: Toward dual-functional wireless networks for 6G and beyond[J].
\textit{IEEE Journal on Selected Areas in Communications}, 2022, 40(6): 1728--1767.

\bibitem{ref_isar_chen}
Chen V C, Martorella M.
\textit{Inverse Synthetic Aperture Radar Imaging: Principles, Algorithms and Applications}[M].
Raleigh, NC: SciTech Publishing, 2014.

\bibitem{ref_radar_signal}
Richards M A.
\textit{Fundamentals of Radar Signal Processing}[M].
2nd ed. New York: McGraw-Hill, 2014.

\bibitem{ref_isar_autofocus}
Wang J, Liu X, Zhou Z.
Minimum-entropy phase adjustment for ISAR[J].
\textit{IEE Proceedings -- Radar, Sonar and Navigation}, 2004, 151(4): 203--209.

\bibitem{ref_pso}
Kennedy J, Eberhart R.
Particle swarm optimization[C]//
\textit{Proceedings of ICNN'95 -- International Conference on Neural Networks}.
Perth, Australia: IEEE, 1995: 1942--1948.

\bibitem{ref_kalman}
Welch G, Bishop G.
An introduction to the Kalman filter[R].
University of North Carolina at Chapel Hill, TR 95-041, 2006.

\bibitem{ref_tdoa}
Chan Y T, Ho K C.
A simple and efficient estimator for hyperbolic location[J].
\textit{IEEE Transactions on Signal Processing}, 1994, 42(8): 1905--1915.

\bibitem{ref_hog}
Dalal N, Triggs B.
Histograms of oriented gradients for human detection[C]//
\textit{IEEE Conference on Computer Vision and Pattern Recognition (CVPR)}.
San Diego, USA: IEEE, 2005: 886--893.

\end{thebibliography}

\vspace{1.5cm}
\begin{center}
  \color{gray}\small
  \href{https://github.com/A1RER/Radar_Proj}{github.com/A1RER/Radar\_Proj}
  \quad$\cdot$\quad MIT License
\end{center}

\end{document}
